\noindent\textbf{\sffamily\large Tóm tắt}\\[1pt]
\highlightlabel{Bối cảnh} Sự gia tăng đột biến và có hệ thống của các ấn phẩm bị rút lại trong hệ sinh thái xuất bản toàn cầu giai đoạn gần đây đang đặt ra thách thức nghiêm trọng đối với tính toàn vẹn khoa học. Điều này đòi hỏi các phương pháp tiếp cận định lượng và trắc lượng thư mục toàn diện nhằm bóc tách xu hướng, nguyên nhân cốt lõi và các mô hình vi phạm đạo đức nghiên cứu trong môi trường học thuật hiện đại.\par\vspace{1pt}

\highlightlabel{Phương pháp} Nghiên cứu này đề xuất một khung quy trình phân tích dữ liệu có khả năng tái lập cao, bắt đầu từ việc thu thập đồng bộ dữ liệu thư mục từ Scopus, Web of Science kết hợp với dữ liệu sự kiện từ Retraction Watch Database (RWD) cho giai đoạn 2021 - 2026.\par\vspace{1pt}

\highlightlabel{Kết quả} Phân tích tập dữ liệu gồm 14.250 ấn phẩm bị rút lại cho thấy số lượng sự kiện thu hồi tăng mạnh theo cấp số nhân, đạt đỉnh vào năm 2024. Về nguyên nhân, các vấn đề liên quan đến dữ liệu (data concerns, chiếm 34,2\%) và gian lận có tổ chức từ các xưởng giấy (paper mills, chiếm 28,5\%) là hai tác nhân chủ đạo. Độ trễ thu hồi trung bình được ghi nhận là 14,5 tháng kể từ thời điểm xuất bản chính thức.\par\vspace{1pt}

\highlightlabel{Kết luận} Phương pháp tiếp cận này cung cấp một khung làm việc toàn diện và có khả năng tái lập cao giúp các nhà quản lý, biên tập viên và nhà nghiên cứu nhận diện sớm các dấu hiệu tiêu cực trong xuất bản học thuật, qua đó góp phần nâng cao tính minh bạch và độ tin cậy của văn liệu khoa học toàn cầu.\par\vspace{12pt}

\highlightlabel{Từ khóa} 
Công bố rút lại \keywordseparator
Trắc lượng thư mục \keywordseparator
Đạo đức xuất bản \keywordseparator
Khoa học Thông tin và Thư viện \keywordseparator
Tính toàn vẹn nghiên cứu \keywordseparator
Khung phân tích tái lập